\documentclass[11pt,a4paper]{article}
\usepackage[margin=1in]{geometry}
\usepackage{hyperref}
\hypersetup{hidelinks} % removes the red boxes hyperref draws around footnote/reference links
\usepackage{enumitem}
\usepackage{graphicx}
\usepackage{xcolor}
\usepackage{array}
\usepackage{booktabs}

% -----------------------------------------------------
% bvraghav-tiet-ucs503p-article.sty
% -----------------------------------------------------

% Variable: Lab Instructor
\makeatletter \newcommand{\BvrSetLabInstructor}[1]{%
  \gdef\Bvr@LabInstructor{#1}}
% default
\newcommand{\Bvr@LabInstructor}{Dr. Raghav %
  B. Venkataramaiyer}
% public accessor
\newcommand{\BvrLabInstructorValue}{\Bvr@LabInstructor}
\makeatother

% Add subtitle
\usepackage{titling}
% -- Set title bold
\pretitle{\begin{center}\bfseries\fontsize{18bp}{18bp}\selectfont}
\makeatletter
\newcommand{\BvrSubtitle}[1]{%
  \posttitle{%
    \par\end{center}
    \begin{center}\large#1\end{center}
    \vskip 0.5em}%
}
\makeatother

% Hints in the section title(s)
\newcommand{\BvrHint}[1]{%
  {\normalfont\normalsize\itshape\hfill #1}}

% -----------------------------------------------------

\title{SafeRoute: A Crowdsourced Safety-Aware Navigation Web App}

\BvrSetLabInstructor{Dr. Paramveer Sidhu}

\BvrSubtitle{%
  Project Proposal for UCS503  \\
  Submitted to: \BvrLabInstructorValue \\ \bfseries
  Thapar Institute of Engineering and Technology% 
}

\author{
  \textsc{Dhriti Handa}, COE%
  \thanks{Roll No: \texttt{1024030176}, %
    Email: \texttt{dhanda\_be24@thapar.edu}}
  \and
  \textsc{Disha Kundra}, COE%
  \thanks{Roll No: \texttt{1024030180}, %
    Email: \texttt{ddisha\_be24@thapar.edu}}
}

\date{\today}

\begin{document}
\maketitle

\section{Higher-order goal%
  \BvrHint{\ldots secondary framing}}
Everyone who has ever walked home at night has done the same
mental calculation: which street actually feels safe, not just
which one is shortest. No navigation tool answers that question
today. SafeRoute does. It's a web application -- open a link, no
install, no app store -- that puts a live, crowdsourced safety
score next to every route, so students, women commuting after
work, and night-shift staff can choose the safer path with the
same confidence they already choose the faster one. The bigger
goal is simple: make safety a routing input, the same way Google
Maps made traffic a routing input.

\section{Time-to-value %
\BvrHint{\ldots primary heuristic}}
\begin{itemize}[leftmargin=0.7cm]
    \item We will deliver a visibly working core loop early: the
    web app will show a live map with a seeded safety heatmap and
    a route that visibly differs between ``fastest'' and
    ``safest'' toggles.
    \item We will prioritize a usable, demoable increment over a
    fully general one -- the first version uses manually seeded
    incident data before the crowdsourcing and ML layers are
    live, so the value of the safety toggle is demonstrable from
    the first prototype.
    \item Success in each iteration is defined by what can be
    shown end-to-end in a browser (navigation, reporting,
    heatmap), not by the completeness of any single subsystem.
\end{itemize}

\section{Problem Statement}
Conventional navigation platforms -- Google Maps, Apple Maps --
are excellent at optimizing for distance, travel time, and traffic.
None of them give a user a dedicated, dynamically-computed safety
layer to compare routes on. This gap matters most exactly when
it's most dangerous to ignore: walking at night or through an
unfamiliar area, where lighting and reported incidents can differ
sharply between two streets that are otherwise equally fast.
\begin{itemize}[leftmargin=0.7cm]
    \item \textbf{No safety option, period:} the shortest path may
    run through a poorly-lit or isolated stretch, and neither the
    app nor the user has any way to know that in advance.
    \item \textbf{Safety tools exist, but don't navigate:} apps
    like Safetipin collect useful safety data, but they function as
    planning/audit dashboards for city governments -- not as a
    consumer tool you'd actually open to get walking directions.
    \item \textbf{SOS tools exist, but don't route:} apps like
    bSafe or Noonlight handle emergencies well, but are purely
    reactive -- they say nothing about which street to take in the
    first place, and require an install.
    \item \textbf{No time-of-day awareness:} a street that's fine
    at 6~pm can feel very different at 11~pm; mainstream routing
    tools don't capture this at all.
\end{itemize}

\textbf{Why this matters right now:} street lighting and safety
perception vary block by block across Indian cities and campuses.
This is something we and our classmates navigate every single day
-- it's a real, felt problem, not one we invented to fit a project
brief, and it currently gets solved with word-of-mouth guesswork
instead of any systematic tool.

\section{Existing Solutions and the Gap We Fill}
\begin{itemize}[leftmargin=0.7cm]
    \item \textbf{Google/Apple Maps:} best-in-class routing and
    real-time traffic, but zero safety-awareness -- optimizes
    purely for time and distance.
    \item \textbf{Safetipin:} pioneered urban safety scoring in
    India and works with city partners, but it's a planning tool,
    not a navigation app -- there's no ``give me the safest route
    right now'' experience, and coverage outside partner cities is
    thin.
    \item \textbf{bSafe / Noonlight:} strong SOS and
    live-location-sharing, but mobile-app-only, install-required,
    and entirely reactive -- no routing component at all.
\end{itemize}
SafeRoute is the first to put these three ideas -- live routing,
crowdsourced safety data, and one-tap SOS -- into a single,
install-free web app, with a spatial model that fills in the gaps
where crowd data doesn't exist yet.

\section{Proposed Solution %
\BvrHint{\ldots Software Proposition}}
\subsection{Overview}
SafeRoute is a \textbf{responsive web application} that layers a
dynamically-updated, crowdsourced safety score on top of standard
map routing:
\begin{itemize}[leftmargin=0.7cm]
    \item \textbf{Turn-by-turn navigation:} standard walking
    directions via a self-hosted OSRM engine, rendered on a Mapbox
    GL JS map, right in the browser.
    \item \textbf{Fastest vs.\ Safest toggle:} one click re-ranks
    candidate routes by a safety-weighted cost function instead of
    distance alone.
    \item \textbf{One-click incident reporting:} geotagged reports
    (poor lighting, isolated area, past incident) using the
    browser's Geolocation API -- minimal friction, no login wall
    to just look around.
    \item \textbf{Live safety heatmap:} a color-graded overlay
    built from real reports, historical density, and time of day.
    \item \textbf{Spatial interpolation for sparse areas:} the map
    is never blank -- unreported streets get an estimated score
    based on nearby data.
    \item \textbf{SOS quick-share:} one click generates a live
    location link sent to a pre-set emergency contact.
\end{itemize}

\subsection{Core workflow}
\begin{enumerate}[leftmargin=0.7cm]
    \item User opens the app in any browser, picks a source and
    destination, and taps ``fastest'' or ``safest.''
    \item The backend checks safety scores for candidate routes in
    PostGIS and re-ranks them with the safety-weighted cost
    function.
    \item While reviewing the route, if the user knows of a
    hazard, they log it in two clicks -- location captured
    automatically via the browser.
    \item Nearby safety-score cells update, and the safety
    heatmap refreshes for the next person who opens that area.
\end{enumerate}

\section{Solution Approach %
\BvrHint{\ldots Engineering focus}}
This is an engineering course project, so our effort goes into
geospatial systems integration and routing correctness -- not into
inventing new machine-learning research.

\subsection{Safety-score interpolation %
\BvrHint{\ldots non-research}}
Estimating safety where reports are sparse is heuristic-and-model
assisted, not a research contribution:
\begin{itemize}[leftmargin=0.7cm]
    \item Raw incident reports are aggregated into a PostGIS grid,
    weighted by recency and community upvotes.
    \item Inverse Distance Weighting (IDW) -- a simple,
    interpretable spatial interpolation technique -- fills in
    scores for grid cells with no direct reports, using nearby
    cells and time-of-day as inputs.
    \item Interpolated cells are always shown as lower-confidence
    than directly-reported ones, and the app never hides or blocks
    a route without the user's own choice.
\end{itemize}

\subsection{System architecture %
\BvrHint{\ldots web + geospatial}}
A React web client talks to a Flask backend, which coordinates a
geospatial database and a routing engine:
\begin{itemize}[leftmargin=0.7cm]
    \item \textbf{Web client:} React + Vite, with Mapbox GL JS for
    the map, incident reporting, the heatmap layer, and SOS --
    configured as an installable PWA.
    \item \textbf{Backend API:} Flask REST API handling auth,
    incident reports, safety-score queries, and safety-weighted
    routing requests.
    \item \textbf{Geospatial database:} PostgreSQL with PostGIS,
    storing grid-based safety scores and reports with spatial
    (GiST) indexing for fast radius/route queries.
    \item \textbf{Routing engine:} a self-hosted OSRM instance;
    the safety layer is applied either by modifying OSRM's cost
    profile, or -- as a fallback if that proves too complex --
    by re-ranking OSRM's candidate routes using an externally
    computed safety score.
    \item \textbf{Realtime layer:} Firebase Auth for login and
    Firebase Cloud Messaging / Web Push for low-safety-zone alerts.
    \item \textbf{Batch interpolation job:} a periodic job
    recomputes the safety grid from new reports using the IDW
    model, keeping live queries fast.
\end{itemize}

\subsection{Tech stack at a glance}
\begin{center}
\begin{tabular}{@{}p{3.2cm}p{10cm}@{}}
\toprule
\textbf{Layer} & \textbf{Tools} \\
\midrule
Frontend & React, Vite, JavaScript, HTML/CSS, Mapbox GL JS \\
Backend & Python, Flask, REST APIs \\
Database & PostgreSQL + PostGIS \\
Routing & OSRM (self-hosted), Mapbox Directions \& Geocoding APIs \\
Auth \& Notifications & Firebase Authentication, Firebase Cloud Messaging / Web Push \\
AI/ML \& Data & Python, NumPy, Pandas, scikit-learn / SciPy (IDW interpolation) \\
Dev Tools & Git, GitHub, VS Code, Postman \\
Deployment & Vercel/Netlify (frontend), Render/Railway (backend + OSRM), Supabase/Neon (PostGIS) \\
\bottomrule
\end{tabular}
\end{center}
Every layer uses a mature, well-documented tool -- we spend our
time on routing correctness and geospatial performance, not on
reinventing infrastructure that's already solved.

\subsection{CI/CD and fast, safe delivery}
\begin{itemize}[leftmargin=0.7cm]
    \item Every push runs backend tests (pytest) and frontend
    tests (Jest + React Testing Library) via GitHub Actions.
    \item Each milestone gets a repeatable frontend deploy
    (Vercel/Netlify) and staging backend deploy ahead of the demo.
    \item New features -- reporting, heatmap, safest-route toggle,
    SOS -- ship incrementally without breaking core navigation.
\end{itemize}

\section{Evaluation Criterion %
\BvrHint{\ldots measurable and attributable}}
We won't just claim SafeRoute works -- we'll measure it.

\subsection{Primary evaluation metric}
\textbf{Route divergence validity:} does the ``safest route''
actually differ from the ``fastest route'' when safety scores
along the corridor are clearly different?
\begin{itemize}[leftmargin=0.7cm]
    \item Target: on a seeded test map of our campus/demo area,
    safest and fastest routes visibly diverge in at least 80\% of
    test cases containing a deliberately low-safety segment.
    \item Measurement: compare route geometry from the
    safety-weighted routing function against the plain,
    unweighted OSRM route for identical start/end points.
\end{itemize}

\subsection{Secondary evaluation metrics}
\begin{itemize}[leftmargin=0.7cm]
    \item \textbf{Report-to-heatmap latency:} time between a
    submitted report and the heatmap reflecting it (target: under
    10 seconds).
    \item \textbf{Interpolation sanity check:} spot-check
    IDW-estimated scores against a small hand-labeled set of
    ``ground truth'' cells.
    \item \textbf{Query performance:} p95 latency for
    radius/route safety queries under simulated load.
    \item \textbf{Reliability:} uptime of the staging/production
    backend and OSRM instance during the demo window.
\end{itemize}

\subsection{Pilot validation plan %
\BvrHint{\ldots high-level}}
\begin{itemize}[leftmargin=0.7cm]
    \item Get 10--20 peers on campus to seed the first round of
    reports and try the reporting flow.
    \item Run a live walkthrough: show fastest and safest routes
    diverging, submit a live report on stage, and watch the
    heatmap update in real time.
\end{itemize}

\section{Scalability %
\BvrHint{\ldots with foresight}}
The system needs to hold up both in theory and in practice.
\begin{enumerate}[leftmargin=0.7cm]
    \item \textbf{Scales in theory:} the web client is fully
    decoupled from the backend and routing engine; PostGIS spatial
    indexing (GiST) replaces full-table scans; the backend API is
    stateless, so it can scale horizontally behind a load balancer.
    \item \textbf{Deployable in practice:} React frontend on
    Vercel/Netlify (static hosting + CDN); Flask backend and OSRM
    on Render/Railway; managed PostGIS-enabled Postgres on
    Supabase or Neon; a GitHub Actions CI/CD pipeline for staging
    and production.
\end{enumerate}
Beyond a single campus, city-scale coverage would need bigger OSRM
map data and query infrastructure -- a deliberate, honest
``future scope'' item, not something we're pretending is already
solved.

\section{Engine availability heuristic}
Nothing here is built from scratch -- every piece rests on a
mature, well-supported tool:
\begin{itemize}[leftmargin=0.7cm]
    \item Web framework: React + Vite, for a fast, map-heavy
    single-page app.
    \item Routing: self-hosted OSRM with a modified cost profile
    (or route re-ranking as fallback), plus Mapbox for map tiles
    and geocoding.
    \item Geospatial storage: PostgreSQL + PostGIS with GiST
    indexing.
    \item Backend: Flask for the REST API.
    \item Auth and push: Firebase Authentication and Web Push /
    Firebase Cloud Messaging.
    \item Interpolation: scikit-learn / SciPy for a lightweight
    IDW implementation, not a custom research model.
\end{itemize}

\section{Project Scope and Deliverables %
\BvrHint{\ldots iteration-friendly}}
\subsection{Initial deliverable %
\BvrHint{\ldots first iteration / prototype}}
\begin{itemize}[leftmargin=0.7cm]
    \item Proposal finalization, wireframes, repository and
    architecture setup.
    \item Basic map and turn-by-turn navigation integration
    (Mapbox/OSRM) in the web app.
    \item Incident reporting feature and PostGIS schema for
    reports and the safety-score grid.
    \item \textbf{Prototype:} web app shows a map with a
    manually-seeded safety heatmap and a route that visibly
    differs between ``fastest'' and ``safest'' toggles, for our
    own campus/city.
\end{itemize}

\subsection{Subsequent deliverables}
\begin{itemize}[leftmargin=0.7cm]
    \item Dynamic safety-score computation from live reports and
    the safest-route toggle wired to the real cost function.
    \item Progress report; IDW-based interpolation model for
    sparse-data areas.
    \item SOS/emergency quick-share feature, push notifications
    for low-safety zones, and UX polish.
    \item Load testing of geospatial queries and bug fixes.
    \item Deployment of the backend, OSRM instance, and web app.
    \item Final presentation preparation and live demo
    rehearsal.
\end{itemize}

\section{Risks and Mitigations}
\begin{itemize}[leftmargin=0.7cm]
    \item \textbf{Cold start (no data for new areas):} manually
    seed safety scores for the campus/demo area up front; mark
    interpolated cells as lower-confidence.
    \item \textbf{Fake or malicious reports:} rate-limit reports
    per authenticated user, require a minimum upvote count before
    a report meaningfully moves the safety score.
    \item \textbf{OSRM cost-function complexity:} budget explicit
    research time in Week 3--4; fallback is re-ranking OSRM's
    candidate routes by an externally computed safety score,
    rather than editing OSRM's Lua profile internally.
    \item \textbf{Browser geolocation limits:} permission prompts
    and desktop IP-based accuracy are worse than mobile GPS --
    mitigated with manual address entry as a graceful fallback.
    \item \textbf{Slow queries at scale:} PostGIS GiST indexes from
    day one, load-tested in Week~13, well before deployment.
    \item \textbf{Two people, three domains (frontend, geospatial
    backend, ML):} strict scope discipline -- PWA polish and push
    notifications are stretch goals, and the Week~7 prototype
    scope is protected above everything else.
\end{itemize}

\section{Summary}
This project proposes a deployable, web-based, safety-aware
navigation application that layers crowdsourced and interpolated
safety scores on top of standard shortest-path routing. It meets
the course criteria by emphasizing time-to-value through a Week-7
demoable prototype, implementing CI/CD to support frequent safe
releases, and using measurable evaluation criteria (route
divergence validity, report-to-heatmap latency, and geospatial
query performance). It remains focused on systems-engineering
depth -- full-stack web development, geospatial databases, and
applied spatial interpolation -- rather than research-driven
algorithmic novelty.

\end{document}
